\documentclass[10pt,twocolumn]{article}
\usepackage[T1]{fontenc}
\usepackage[utf8]{inputenc}
\usepackage{lmodern}
\usepackage[margin=1.8cm,top=2.4cm,bottom=2cm,columnsep=0.6cm]{geometry}
\usepackage{fancyhdr}
\usepackage{titlesec}
\usepackage{booktabs}
\usepackage{amsmath,amssymb,amsfonts}
\usepackage{microtype}
\usepackage{xcolor}
\usepackage{mdframed}
\usepackage{parskip}
\usepackage{enumitem}
\usepackage{hyperref}

\definecolor{rulered}{RGB}{180,30,30}
\definecolor{boxgray}{RGB}{245,245,245}
\definecolor{keywordgray}{RGB}{100,100,100}

\pagestyle{fancy}
\fancyhf{}
\fancyhead[L]{\textsc{\small Claude Mag}}
\fancyhead[C]{\small\textit{Machine Learning Research Digest}}
\fancyhead[R]{\small 2026-11-11}
\fancyfoot[C]{\small\thepage}
\renewcommand{\headrulewidth}{0.8pt}

\titleformat{\section}{\normalsize\bfseries\color{rulered}}{}{0pt}{}%
  [\vspace{-4pt}\color{rulered}\rule{\columnwidth}{0.4pt}]
\titlespacing{\section}{0pt}{8pt}{4pt}

\hypersetup{colorlinks=true,urlcolor=rulered,linkcolor=black}

\begin{document}
\setlength{\parindent}{0pt}
\setlength{\parskip}{4pt}

{\small\color{keywordgray}\textbf{Keywords:}
  graph neural networks \textbullet{}
  Muon optimizer \textbullet{}
  second-order methods \textbullet{}
  graph-aware preconditioning}\par\vspace{4pt}

{\LARGE\bfseries\raggedright
  Module-Aware Optimization for Graph Neural Networks}\par\vspace{4pt}

{\large\itshape\raggedright
  GraphMuon Adapts the Muon Optimizer to Graph-Propagated Input Covariance,
  Achieving 2.13--7.38$\times$ Wall-Clock Speedups Over AdamW}\par\vspace{6pt}

{\small\textbf{By} Guy Hadad, Haggai Roitman, Moshe Eliasof, Bracha Shapira}\par\vspace{2pt}

{\small\color{keywordgray}arXiv:2611.08379 \textbullet{}
  cs.LG / cs.AI \textbullet{} NeurIPS 2026}
\par\vspace{2pt}\noindent{\color{rulered}\rule{\columnwidth}{0.6pt}}\vspace{6pt}

The Muon optimizer orthogonalizes gradient matrices to precondition updates,
but its implicit geometry ignores graph propagation, which systematically
distorts layer-input covariance in GNNs. GraphMuon corrects this by
orthogonalizing in the graph-propagated input metric, cutting
time-to-target by up to 7.38$\times$ over AdamW on 15 benchmark
configurations.

\begin{mdframed}[backgroundcolor=boxgray,linecolor=rulered,linewidth=1pt,
    innerleftmargin=6pt,innerrightmargin=6pt,
    innertopmargin=6pt,innerbottommargin=6pt]
\textbf{\small AT A GLANCE}\par\vspace{4pt}
\begin{description}[leftmargin=0pt,labelsep=4pt,itemsep=2pt,parsep=0pt,
    font=\small\bfseries]
  \item[Problem:] Muon's matrix update assumes isotropic input covariance, but
    graph propagation distorts covariance, amplifying updates and destabilizing
    GNN training.
  \item[Why it matters:] GNNs are increasingly used for large-scale scientific
    and industrial graph tasks; faster, stable optimization directly reduces
    training cost.
  \item[Method:] Replace Muon's Euclidean metric with the graph-propagated input
    covariance; orthogonalize and normalize step sizes in this graph-aware metric.
  \item[Result:] 2.13--7.38$\times$ wall-clock speedup over AdamW on 13/15
    dataset-architecture configs; vanilla Muon diverges on 4/15.
  \item[Evidence:] 15 configurations across node classification, inductive graph
    classification, and graph regression benchmarks.
\end{description}
\end{mdframed}

\section{The Big Picture}

The Muon optimizer has emerged as a competitive alternative to AdamW for
large neural network training. Its core mechanism is matrix orthogonalization:
rather than scaling the gradient element-wise (as Adam does), Muon computes the
polar decomposition of the gradient matrix and uses the orthogonal factor as the
update. This is equivalent to preconditioned gradient descent with preconditioner
$(GG^T)^{-1/2}$, and it implicitly implements a form of second-order optimization
using only the gradient's own geometry.

Muon was developed and evaluated on transformer architectures, where the input
to each linear layer has approximately isotropic covariance (thanks to layer
normalization before attention and MLP blocks). This isotropy is critical: it
ensures that orthogonalizing the gradient matrix produces an update whose
magnitude is well-calibrated relative to the effect on the layer's output.

Graph neural networks break this assumption. In a GNN layer, the input to each
linear transformation is the aggregated neighborhood feature $H = \tilde{A}X$,
where $\tilde{A}$ is the normalized adjacency operator. Graph propagation
concentrates variance in directions corresponding to low-frequency graph
components (smooth node features) and suppresses high-frequency directions.
Applying Muon's Euclidean orthogonalization to the gradient of $WH^T$ ignores
this distortion, causing the update to be under-calibrated in smooth directions
and over-calibrated in rough directions --- destabilizing training in poorly
conditioned graphs.

\section{Why Existing Approaches Fall Short}

\textbf{AdamW} applies element-wise adaptive scaling, which is permutation-
invariant and does not exploit the matrix structure of weight gradients.
It is robust but slow to converge on tasks where curvature is anisotropic
along structured directions.

\textbf{Vanilla Muon} captures matrix structure but uses a Euclidean metric
that is mismatched to graph-propagated inputs. On graphs with high spectral
anisotropy (e.g., large homophilic graphs with strong community structure),
the covariance of $H$ is far from isotropic, and Muon's updates amplify
variance along over-represented directions, causing divergence.

\textbf{SOAP and Shampoo} are full Kronecker-factored second-order methods that
in principle handle anisotropic curvature, but their cost per step is cubic in
the layer dimension, making them impractical for large GNN layers.

GraphMuon achieves the efficiency of Muon (one Newton-Schulz iteration per
step, $O(d^2)$ cost) while correcting for the GNN-specific covariance structure.

\section{Core Mechanism}

For a GNN layer with propagated input $H = \tilde{A}X \in \mathbb{R}^{n \times d}$
and weight matrix $W \in \mathbb{R}^{d' \times d}$, the layer output is $Z = HW^T$.
The gradient of the loss with respect to $W$ is:
\[
  G = \frac{\partial \mathcal{L}}{\partial W} = \left(\frac{\partial \mathcal{L}}{\partial Z}\right)^T H.
\]
The empirical input covariance is $\Sigma = H^T H / n \in \mathbb{R}^{d\times d}$.

Muon computes the polar factorization $G = UP$ (with $U$ orthogonal and $P$
positive semi-definite) and uses $U$ as the update. This is the orthogonalization
in the standard Frobenius metric $\langle A, B \rangle_F = \text{tr}(A^T B)$.

GraphMuon instead orthogonalizes in the $\Sigma$-weighted metric:
\[
  \langle A, B \rangle_\Sigma = \text{tr}(A \Sigma B^T).
\]
The $\Sigma$-orthogonal factor of $G$ is obtained by computing the polar
factorization of $G\Sigma^{1/2}$:
\[
  G\Sigma^{1/2} = \hat{U}\hat{P}, \quad
  U_\Sigma = \hat{U}\Sigma^{-1/2},
\]
where $\hat{U}$ is orthogonal and $\hat{P}$ is PSD. The step-size normalization
replaces the spectral-norm factor used in Muon with:
\[
  \eta_\Sigma = \frac{\eta}{\|U_\Sigma H^T\|_F},
\]
ensuring the effective learning rate measures actual change in layer output.

\section{Technical Formulation}

Let $\Sigma_t$ be an exponential moving average of the input covariance:
\[
  \Sigma_t = \beta \Sigma_{t-1} + (1-\beta)\, \frac{H_t^T H_t}{n_t},
\]
with $\beta = 0.95$. The GraphMuon update at step $t$ is:
\[
  W_{t+1} = W_t - \eta_\Sigma\, U_{\Sigma_t},
\]
where $U_{\Sigma_t}$ is computed via two Newton-Schulz iterations applied to
$G_t \Sigma_t^{1/2}$:
\[
  B_0 = G_t \Sigma_t^{1/2},\quad
  B_{k+1} = \tfrac{3}{2}B_k - \tfrac{1}{2}B_k B_k^T B_k,\quad
  U_{\Sigma_t} = B_2 \Sigma_t^{-1/2}.
\]
The $\Sigma_t^{1/2}$ and $\Sigma_t^{-1/2}$ factors are computed via eigendecomposition
once per update, sharing the cost across all gradient steps in a mini-batch.
Per-step overhead over vanilla Muon is dominated by the matrix square roots
and is bounded by $O(d^3 / n_\text{steps})$ amortized.

\section{Learning or Inference Procedure}

GraphMuon is a drop-in replacement for AdamW or Muon in any GNN training
pipeline. Required inputs are:
\begin{itemize}[itemsep=2pt,parsep=0pt]
  \item Gradient matrices $G_t$ for each GNN weight matrix, computed by
    standard backpropagation.
  \item Propagated activations $H_t$ for each layer, cached from the
    forward pass.
  \item An EMA coefficient $\beta$ (default 0.95) and two Newton-Schulz
    iterations (sufficient for convergence at standard scales).
\end{itemize}
No separate profiling or warm-up phase is required; the EMA initializes from the
first mini-batch and stabilizes within a few hundred steps.

\section{What the Guarantee Says}

No formal convergence rate is derived. The empirical claim is: on 13 of 15
evaluated dataset-architecture configurations, GraphMuon reaches a fixed
performance target faster (in wall-clock time) than AdamW. On the remaining 2,
it is within 5\% of AdamW's time-to-target. Vanilla Muon fails to converge on
4 of the 15 configurations; GraphMuon converges on all 15.

\section{Experimental Findings}

\begin{itemize}[itemsep=2pt,parsep=0pt]
  \item \textbf{Time-to-target:} 2.13--7.38$\times$ faster than AdamW
    (13/15 configs); speedups are largest on dense homophilic graphs.
  \item \textbf{Peak accuracy:} Marginally higher than AdamW in 9/15 configs;
    never more than 0.3\% lower in any configuration.
  \item \textbf{Vanilla Muon failures:} Muon diverges on 4/15 configs
    (ogbn-arxiv with GraphSAGE, ZINC with GNN+, and 2 TUDatasets); GraphMuon
    converges on all.
  \item \textbf{Overhead:} Per-step overhead over AdamW is 8--12\% in
    wall-clock time; the speedup from fewer steps more than compensates.
\end{itemize}

\section{Ablations and Interpretation}

\textbf{EMA vs.\ online covariance:} Using a fresh covariance estimate per step
achieves the same quality but at 2.3$\times$ the per-step cost; EMA captures
the same effective geometry at lower cost.

\textbf{Newton-Schulz iterations:} One iteration is insufficient for high spectral
ratio $\Sigma$; two iterations yield near-exact $\Sigma$-orthogonalization
in all evaluated cases.

\textbf{Step-size normalization:} Replacing the output-space normalization
with the standard Muon spectral-norm normalization degrades time-to-target by
18\% on average, confirming the importance of measuring step size in output space.

\textbf{Covariance estimation noise:} On sparse graphs (fewer than 100 nodes
per batch), the EMA estimate is noisy; adding a 5\% identity regularizer
to $\Sigma_t$ recovers the full speedup without hurting peak accuracy.

\section*{Reference}

Guy Hadad, Haggai Roitman, Moshe Eliasof, Bracha Shapira.
\textbf{Module-Aware Optimization for Graph Neural Networks.}
arXiv:2611.08379. NeurIPS 2026.
\url{https://arxiv.org/abs/2611.08379}

\end{document}
